\begin{proof}
    Let $x_{1}, x_{2} \in E$ be given.
    \begin{enumerate}[label=\textbf{(\Alph*)}]
        \item If $f$ is injective, then
        \[
        \begin{array}{c}
            f|_{E}^{F}(x_{1}) = f|_{E}^{F}(x_{2}) \\
            \Downarrow \\
            f(x_{1}) = f(x_{2}) \\
            \Downarrow \\
            x_{1} = x_{2}
        \end{array}
        \]
        \item If $f(E) = F$, then
        \[
        \begin{array}{c}
            \forall y \in F, \exists x \in E, f(x) = y \\
            \Downarrow \\
            \forall y \in F, \exists x \in E, f|_{E}^{F}(x) = y
        \end{array}
        \]
        \item follows from \textbf{(A)} and \textbf{(B)}
    \end{enumerate}
\end{proof}

\bigskip

\begin{definition} \label{an:def:composition-function}
    Let $f: X \to Y$ and $g: Y \to Z$ be given. The \textbf{composition} $g \circ f: X \to Z$ of $g$ with $f$ is defined as
    \[
    \forall x \in X, (g \circ f)(x) = g(f(x))
    \]
\end{definition}

\bigskip

\begin{theorem} \label{an:thm:properties-of-composition-functions}
    Let $f: X \to Y$ and $g: Y \to Z$ be given. Then the following properties hold:
    \begin{enumerate}[label=\textbf{(\Alph*)}]
        \item $f$ and $g$ are injective $\implies g \circ f$ is injective
        \item $f$ and $g$ are surjective $\implies g \circ f$ is surjective
        \item $f$ and $g$ are bijective $\implies g \circ f$ is bijective
    \end{enumerate}
\end{theorem}

\begin{proof}
    Let $x_{1}, x_{2} \in X$ be given.
    \begin{enumerate}[label=\textbf{(\Alph*)}]
        \item If $f$ and $g$ are injective, then
        \[
        \begin{array}{c}
            (g \circ f)(x_{1}) = (g \circ f)(x_{2}) \\
            \Downarrow \\
            g(f(x_{1})) = g(f(x_{2})) \\
            \Downarrow \\
            f(x_{1}) = f(x_{2}) \\
            \Downarrow \\
            x_{1} = x_{2}
        \end{array}
        \]
        \item If $f$ and $g$ are surjective, then
        \[
        \forall z \in Z, \exists y \in Y, g(y) = z
        \]
        Also,
        \[
        \forall y \in Y, \exists x \in X, f(x) = y
        \]
        Therefore,
        \[
        \forall z \in Z, \exists x \in X, (g \circ f)(x) = z
        \]
        \item follows from \textbf{(A)} and \textbf{(B)}
    \end{enumerate}
\end{proof}

\begin{theorem} \label{an:thm:bijection-equivalence-relation}
    Let $S$ be a set of sets. Define
    \[
    \forall X, Y \in S, X \sim Y \iff \exists f: X \xrightarrow{\sim} Y
    \]
    Then $\sim$ is an equivalence relation on $S$.
\end{theorem}

\begin{proof}
    Let $X, Y, Z \in S$ be given.
    \begin{itemize}
        \item \textbf{Reflexivity:}
        \[
        \begin{array}{c}
            \forall X \in S, \exists f: X \xrightarrow{\sim} X, f(x) = x \\
            \Downarrow \\
            \forall X \in S, X \sim X
        \end{array}
        \]
        \item \textbf{Symmetry:} If $X \sim Y$, by \Cref{an:thm:inverse-function-is-bijective},
        \[
        \begin{array}{c}
            \exists f: X \xrightarrow{\sim} Y \\
            \Downarrow \\
            \exists f^{-1}: Y \xrightarrow{\sim} X
        \end{array}
        \]
        Therefore,
        \[
        \forall X, Y \in S, X \sim Y \implies Y \sim X
        \]
        \item \textbf{Transitivity:} If $X \sim Y$ and $Y \sim Z$, by \Cref{an:thm:inverse-function-is-bijective},
        \[
        \begin{array}{c}
            \exists f: X \xrightarrow{\sim} Y \land \exists g: Y \xrightarrow{\sim} Z \\
            \Downarrow \\
            \exists g \circ f: X \xrightarrow{\sim} Z
        \end{array}
        \]
        Therefore,
        \[
        \forall X, Y, Z \in S, X \sim Y \land Y \sim Z \implies X \sim Z
        \]
    \end{itemize}
\end{proof}
